\documentclass[12pt,letterpaper]{article}

% --- Configuración de Idioma y Tipografía Industrial ---
\usepackage[T1]{fontenc} % Manejo correcto de fuentes con tildes y eñes
\usepackage[utf8]{inputenc}
\usepackage[spanish]{babel}
\usepackage{courier} 
\renewcommand{\familydefault}{\ttdefault} % Todo el documento en Courier

% --- Alternativa robusta para evitar desbordamiento de márgenes ---
\usepackage{ragged2e} % Paquete real y correcto para alineación izquierda limpia

% --- Márgenes Estándar de la Industria (Hollywood/Celtx) ---
\usepackage{geometry}
\geometry{
    letterpaper,
    left=1.5in,   % Margen amplio para encuadernado
    right=1.0in,
    top=1.0in,
    bottom=1.0in
}

% --- Espaciado ---
\usepackage{setspace}
\setstretch{1.15} 
\parindent 0in    

% --- Contador Automático de Escenas ---
\newcounter{numescena}
\setcounter{numescena}{0}

% --- CLON DE ENCABEZADO CELTX ---
\newcommand{\escena}[3]{%
    \refstepcounter{numescena}%
    \vspace{1.5em} 
    \noindent\textbf{\arabic{numescena}. #1. #2 - #3}%
    \vspace{0.8em}\par 
}

% --- FORMATO DE DIÁLOGO ESTILO CELTX ---
\newcommand{\personaje}[1]{%
    \vspace{0.6em}
    \noindent\hspace{2.2in}\textbf{#1}\par
}

% Formato para la Acotación / Parentrencial
\newcommand{\parentesis}[1]{%
    \noindent\hspace{1.6in}(#1)\par
    \vspace{-0.3em}
}

% Formato para el Bloque de Diálogo
\newcommand{\dialogo}[1]{%
    \noindent\hspace{1.0in}\parbox{3.5in}{\RaggedRight #1}\par
    \vspace{0.6em}
}

% --- PORTADA ESTILO CELTX ---
\begin{document}
\RaggedRight % Fuerza la alineación de guion a la izquierda desde el inicio

\thispagestyle{empty}
\vspace*{2.5in}
\begin{center}
    {\Huge \textbf{[TÍTULO DE TU OBRA]}}\\[1cm]
    {\large Escaleta / Guion de Producción Final}\\[2cm]
    Escrito por\\
    \textbf{[Tu Nombre o Seudónimo]}
\end{center}

\vfill
\noindent Contacto:\\\relax
[Tu Email]\\\relax
[Tu Teléfono]\hfill \today
\newpage

% --- INICIO DEL LIBRETO ---
\pagestyle{plain}
\setcounter{page}{1}

\escena{INT}{HABITACIÓN DE CARLOS}{DÍA}
Carlos (25) duerme plácidamente bocarriba, tapado hasta la nariz. El despertador digital de la mesita de noche marca las 08:55 AM y comienza a emitir un pitido ensordecedor.

Carlos lanza un manotazo a ciegas, tirando el aparato al suelo. El cristal se rompe. Abre un ojo, mira la hora entre las grietas y salta de la cama de un brinco.

\personaje{CARLOS}
\parentesis{asustado}
¡No, no, no! ¡Llego tardísimo!

Instintivamente se tropieza con las sábanas y cae de rodillas al suelo.

\escena{EXT}{CALLES DE LA CIUDAD}{DÍA}
Carlos corre desesperado por la acera esquivando peatones. Lleva la camisa mal abotonada y los cordones de los zapatos arrastrando por el suelo.

El tráfico en la avenida principal está completamente colapsado. Las bocinas retumban. A lo lejos, el autobús que necesita cierra las puertas y arranca.

\personaje{CARLOS}
\dialogo{¡Espere! ¡Por favor, espere!}

Carlos acelera el paso con el rostro desencajado, pero el vehículo se pierde entre el humo del escape.

\escena{INT}{EDIFICIO CORPORATIVO - VESTÍBULO}{DÍA}
Las puertas giratorias de cristal se mueven a toda velocidad. Carlos entra derrapando en el suelo de mármol pulido, empapado en sudor. Se detiene en seco frente al mostrador de recepción.

\end{document}